%=====================================================================
%  导言区（中文精译 + 学术审校版）
%  原文：P. Etingof, "Mathematical ideas and notions of quantum field theory",
%        arXiv:2409.03117v1（MIT 数学系研究生课程讲义，2002 年与 2023 年讲授）
%  编译：XeLaTeX 内核（本目录内用 tectonic -X compile main.tex）
%  说明：本文件只含导言区，由 main.tex 通过 %=====================================================================
%  导言区（中文精译 + 学术审校版）
%  原文：P. Etingof, "Mathematical ideas and notions of quantum field theory",
%        arXiv:2409.03117v1（MIT 数学系研究生课程讲义，2002 年与 2023 年讲授）
%  编译：XeLaTeX 内核（本目录内用 tectonic -X compile main.tex）
%  说明：本文件只含导言区，由 main.tex 通过 %=====================================================================
%  导言区（中文精译 + 学术审校版）
%  原文：P. Etingof, "Mathematical ideas and notions of quantum field theory",
%        arXiv:2409.03117v1（MIT 数学系研究生课程讲义，2002 年与 2023 年讲授）
%  编译：XeLaTeX 内核（本目录内用 tectonic -X compile main.tex）
%  说明：本文件只含导言区，由 main.tex 通过 %=====================================================================
%  导言区（中文精译 + 学术审校版）
%  原文：P. Etingof, "Mathematical ideas and notions of quantum field theory",
%        arXiv:2409.03117v1（MIT 数学系研究生课程讲义，2002 年与 2023 年讲授）
%  编译：XeLaTeX 内核（本目录内用 tectonic -X compile main.tex）
%  说明：本文件只含导言区，由 main.tex 通过 \input{preamble} 读入。
%        宏定义 macros.tex 逐字取自原文导言区，保证记号与原文完全一致。
%=====================================================================

%---------------------------------------------------------------------
% 一、页面与版式（A4，便于打印与屏幕阅读）
%---------------------------------------------------------------------
\usepackage[a4paper,left=2.6cm,right=2.6cm,top=2.6cm,bottom=2.6cm]{geometry}
\linespread{1.12}
\setlength{\emergencystretch}{3em}
\tolerance=2000
\setlength{\parindent}{2em}
\setlength{\parskip}{0.2ex}

%---------------------------------------------------------------------
% 二、中文支持（必须使用 XeLaTeX 内核）
%---------------------------------------------------------------------
\usepackage[UTF8,fontset=fandol]{ctex}
% 章标题格式：第 N 章  <标题>（居中）；小节标题保持 article 默认样式
\usepackage{titlesec}
\titleformat{\section}{\Large\bfseries\centering}{第\,\thesection\,章}{1em}{}
\titlespacing*{\section}{0pt}{3.2ex plus 1ex minus .2ex}{1.8ex plus .2ex}
\titleformat{\subsection}{\large\bfseries}{\thesubsection}{0.8em}{}
\titlespacing*{\subsection}{0pt}{2.4ex plus .6ex minus .2ex}{1ex plus .2ex}

%---------------------------------------------------------------------
% 三、数学宏包
%---------------------------------------------------------------------
\usepackage{amsmath}
\usepackage{amsthm}
\usepackage{amssymb}
\usepackage{amsfonts}
\usepackage{latexsym}
\usepackage{mathtools}

%---------------------------------------------------------------------
% 四、图形、颜色、表格、链接
%---------------------------------------------------------------------
\usepackage{graphicx}
\usepackage{xcolor}
\usepackage{booktabs}
\usepackage{array}
\usepackage{enumitem}
\setlist{nosep,leftmargin=*}

%---------------------------------------------------------------------
% 五、定理类环境（名称汉化，编号方式与原文完全一致：共用定理计数器、按章编号）
%---------------------------------------------------------------------
\theoremstyle{plain}
\newtheorem{theorem}{定理}[section]
\newtheorem{lemma}[theorem]{引理}
\newtheorem{exercise}[theorem]{习题}
\newtheorem{proposition}[theorem]{命题}
\newtheorem{corollary}[theorem]{推论}

\theoremstyle{definition}
\newtheorem{definition}[theorem]{定义}
\newtheorem{example}[theorem]{例}
\newtheorem{conjecture}[theorem]{猜想}

\theoremstyle{remark}
\newtheorem{remark}[theorem]{注}
\newtheorem{note}[theorem]{说明}

% 证明环境
\renewcommand{\proofname}{证明}

% 公式按章编号：(1.1), (1.2), ...
\numberwithin{equation}{section}

%---------------------------------------------------------------------
% 六、原文自带的全部记号缩写（逐字照搬，切勿改动）
%---------------------------------------------------------------------
\input{macros}

%---------------------------------------------------------------------
% 七、图表位置标注：在中文图题后自动注明原文页码与图号
%     用法：\caption{中文图题\figorig{原文页码}{原文图号}}
%---------------------------------------------------------------------
\newcommand{\figorig}[2]{\\[1pt]{\footnotesize\color{gray}（原文：第 #1 页，Figure #2）}}

%---------------------------------------------------------------------
% 八、每章“本节重点”环境（译者为方便阅读所加，非原文内容）
%---------------------------------------------------------------------
\newenvironment{keytakeaways}
  {\par\medskip\noindent
   \begin{minipage}{0.96\textwidth}
   \color{blue!45!black}\rule{\linewidth}{0.8pt}\par\smallskip
   \noindent\textbf{本章要点（Key Takeaways）}\par\smallskip\small}
  {\par\smallskip\noindent\color{blue!45!black}\rule{\linewidth}{0.8pt}\end{minipage}\medskip}

%---------------------------------------------------------------------
% 九、审校批注宏（正文一律用 % [Tutor Note: ...] 注释，此宏仅为备用）
%---------------------------------------------------------------------
\newcommand{\TutorNote}[1]{}

%---------------------------------------------------------------------
% 十、超链接（必须最后加载）
%---------------------------------------------------------------------
\usepackage[unicode,colorlinks=true,
            linkcolor=blue!55!black,
            citecolor=green!45!black,
            urlcolor=blue!60!black,
            pdftitle={量子场论的数学思想与概念},
            pdfauthor={Pavel Etingof（中译者）},
            bookmarksnumbered=true]{hyperref}

 读入。
%        宏定义 macros.tex 逐字取自原文导言区，保证记号与原文完全一致。
%=====================================================================

%---------------------------------------------------------------------
% 一、页面与版式（A4，便于打印与屏幕阅读）
%---------------------------------------------------------------------
\usepackage[a4paper,left=2.6cm,right=2.6cm,top=2.6cm,bottom=2.6cm]{geometry}
\linespread{1.12}
\setlength{\emergencystretch}{3em}
\tolerance=2000
\setlength{\parindent}{2em}
\setlength{\parskip}{0.2ex}

%---------------------------------------------------------------------
% 二、中文支持（必须使用 XeLaTeX 内核）
%---------------------------------------------------------------------
\usepackage[UTF8,fontset=fandol]{ctex}
% 章标题格式：第 N 章  <标题>（居中）；小节标题保持 article 默认样式
\usepackage{titlesec}
\titleformat{\section}{\Large\bfseries\centering}{第\,\thesection\,章}{1em}{}
\titlespacing*{\section}{0pt}{3.2ex plus 1ex minus .2ex}{1.8ex plus .2ex}
\titleformat{\subsection}{\large\bfseries}{\thesubsection}{0.8em}{}
\titlespacing*{\subsection}{0pt}{2.4ex plus .6ex minus .2ex}{1ex plus .2ex}

%---------------------------------------------------------------------
% 三、数学宏包
%---------------------------------------------------------------------
\usepackage{amsmath}
\usepackage{amsthm}
\usepackage{amssymb}
\usepackage{amsfonts}
\usepackage{latexsym}
\usepackage{mathtools}

%---------------------------------------------------------------------
% 四、图形、颜色、表格、链接
%---------------------------------------------------------------------
\usepackage{graphicx}
\usepackage{xcolor}
\usepackage{booktabs}
\usepackage{array}
\usepackage{enumitem}
\setlist{nosep,leftmargin=*}

%---------------------------------------------------------------------
% 五、定理类环境（名称汉化，编号方式与原文完全一致：共用定理计数器、按章编号）
%---------------------------------------------------------------------
\theoremstyle{plain}
\newtheorem{theorem}{定理}[section]
\newtheorem{lemma}[theorem]{引理}
\newtheorem{exercise}[theorem]{习题}
\newtheorem{proposition}[theorem]{命题}
\newtheorem{corollary}[theorem]{推论}

\theoremstyle{definition}
\newtheorem{definition}[theorem]{定义}
\newtheorem{example}[theorem]{例}
\newtheorem{conjecture}[theorem]{猜想}

\theoremstyle{remark}
\newtheorem{remark}[theorem]{注}
\newtheorem{note}[theorem]{说明}

% 证明环境
\renewcommand{\proofname}{证明}

% 公式按章编号：(1.1), (1.2), ...
\numberwithin{equation}{section}

%---------------------------------------------------------------------
% 六、原文自带的全部记号缩写（逐字照搬，切勿改动）
%---------------------------------------------------------------------
% 本文件逐字取自原文导言区（qftlect.tex 第 23--133 行）中的宏定义，
% 保证所有记号、花体字、缩写与原文完全一致。

%%%%%%%%%%%%%%% symbols shorthand %%%%%%%%%%%  
%\renewcommand{\text}[1]{\mbox{{\rm #1}}}
\newcommand{\Tr}{\text{Tr}}%{\text{Tr}\,} 
\newcommand{\id}{\text{id}} 
\newcommand{\Image}{\text{Im}} 
\newcommand{\Ker}{\text{Ker\,}} 
\newcommand{\rank}{\text{rank\,}} 
\newcommand{\Index}{\text{Index}\,} 
\newcommand{\Map}{\text{Map}}
\newcommand{\End}{\text{End}} 
\newcommand{\PGL}{\text{PGL}}
\newcommand{\Hom}{\text{Hom}} 
\newcommand{\Ad}{\text{Ad}}
\newcommand{\Aut}{\text{Aut}}
\newcommand{\Sh}{\text{Sh}}
\newcommand{\Ind}{\text{Ind}}
\newcommand{\Rep}{\text{Rep}}
\newcommand{\supp}{\text{supp}\,} 
\renewcommand{\span}{\text{span}} 
\newcommand{\depth}{\text{depth}} 
\newcommand{\St}{\text{St}}
\newcommand{\diag}{\text{diag}}
\newcommand{\comp}{\text{comp}}
\newcommand{\op}{\text{op}}
\newcommand{\tr}{\text{tr}\,} 

   
\newcommand{\eps}{\varepsilon}
\newcommand{\mR}{{\mathcal R}} 
\newcommand{\bmR}{{\bar \mR}}
\newcommand{\bR}{{\bar R}}   
\newcommand{\bV}{{\bar V}}
\newcommand{\bF}{{\bar F}}
\newcommand{\I}{\mathcal{I}}
\renewcommand{\O}{\mathcal{O}}

\newcommand{\G}{\mathcal{G}}
\newcommand{\E}{\mathcal{E}}

\newcommand{\g}{\mathfrak{g}}
\newcommand{\h}{\mathfrak{h}}
\newcommand{\n}{\mathfrak{n}}
\renewcommand{\b}{\mathfrak{b}}
\newcommand{\wg}{\widehat{\mathfrak{g}}}
\newcommand{\wh}{\widehat{\mathfrak{h}}}
\newcommand{\wn}{\widehat{\mathfrak{n}}}
\newcommand{\wb}{\widehat{\mathfrak{b}}}
\newcommand{\ov}{\overline}

\newcommand{\actl}{\rightharpoonup}
\newcommand{\actr}{\leftharpoonup} 
\newcommand{\lact}{\triangleright} 
\newcommand{\ract}{\triangleleft} 
\newcommand{\la}{\langle\,} 
\newcommand{\ra}{\,\rangle} 
\newcommand{\flip}{\varsigma} 
\newcommand{\nv}{\vec{n}} 
\renewcommand{\star}{{\dag}} 
%\newcommand{\rtimes}{{>\!\!\!\triangleleft}} 
%\newcommand{\ltimes}{{\triangleright\!\!\!<}} 
\newcommand{\smrtimes}{>\!\!\triangleleft} 
\newcommand{\equivalent}{\Leftrightarrow} 
\newcommand{\half}{1/2}%%%%%{\frac{1}{2}}   
\newcommand{\own}{\ni}  
\newcommand{\1}{_{(1)}} 
\newcommand{\2}{_{(2)}} 
\newcommand{\3}{_{(3)}} 
\newcommand{\4}{_{(4)}}
\newcommand{\5}{_{(5)}} 
\newcommand{\6}{_{(6)}}
\renewcommand{\I}{^{(1)}} 
\newcommand{\II}{^{(2)}}   
\newcommand{\tDelta}{\Delta^{op}} 
\newcommand{\tS}{\widetilde{S}} 
\newcommand{\teps}{\widetilde\eps}  
\newcommand{\tI}{_{(\widetilde{1})}} 
\newcommand{\tII}{_{(\widetilde{2})}} 
\newcommand{\tIII}{_{(\widetilde{3})}}
\newcommand{\bTheta}{{\bar \Theta}}
\newcommand{\bs}{\backslash} 


\newcommand{\ot}{\otimes}
\newcommand{\qexp}{{\text{qexp}}}
\newcommand{\ben}{\begin{enumerate}}
\newcommand{\een}{\end{enumerate}}
\newcommand{\ad}{{\text{Ad}}}
\newcommand{\Rad}{{\text{Rad}}}
\newcommand{\Vect}{{\text{Vect}}}
\newcommand{\Lie}{{\text{Lie}}}
\newcommand{\Cl}{{\text{Cl}}}
\newcommand{\mC}{{\mathcal C}}
\newcommand{\trace}{{\rm tr}}
\newcommand{\CC}{{\mathbb{C}}}
\newcommand{\mM}{{\mathcal M}}

\hyphenation{se-mi-simple co-se-mi-simple}

%%%%%%% Bonnie added %%%%%%
\renewcommand{\Bbb}{\mathbb}
\renewcommand{\frak}{\mathfrak}
\renewcommand{\bold}{\mathbf}

%%%\newcommand\topfraction{.7}
\renewcommand\topfraction{.6}
%%%\newcommand\bottomfraction{.3}
\renewcommand\bottomfraction{.6}
\numberwithin{equation}{section}
\pagestyle{plain}

\makeindex


%---------------------------------------------------------------------
% 七、图表位置标注：在中文图题后自动注明原文页码与图号
%     用法：\caption{中文图题\figorig{原文页码}{原文图号}}
%---------------------------------------------------------------------
\newcommand{\figorig}[2]{\\[1pt]{\footnotesize\color{gray}（原文：第 #1 页，Figure #2）}}

%---------------------------------------------------------------------
% 八、每章“本节重点”环境（译者为方便阅读所加，非原文内容）
%---------------------------------------------------------------------
\newenvironment{keytakeaways}
  {\par\medskip\noindent
   \begin{minipage}{0.96\textwidth}
   \color{blue!45!black}\rule{\linewidth}{0.8pt}\par\smallskip
   \noindent\textbf{本章要点（Key Takeaways）}\par\smallskip\small}
  {\par\smallskip\noindent\color{blue!45!black}\rule{\linewidth}{0.8pt}\end{minipage}\medskip}

%---------------------------------------------------------------------
% 九、审校批注宏（正文一律用 % [Tutor Note: ...] 注释，此宏仅为备用）
%---------------------------------------------------------------------
\newcommand{\TutorNote}[1]{}

%---------------------------------------------------------------------
% 十、超链接（必须最后加载）
%---------------------------------------------------------------------
\usepackage[unicode,colorlinks=true,
            linkcolor=blue!55!black,
            citecolor=green!45!black,
            urlcolor=blue!60!black,
            pdftitle={量子场论的数学思想与概念},
            pdfauthor={Pavel Etingof（中译者）},
            bookmarksnumbered=true]{hyperref}

 读入。
%        宏定义 macros.tex 逐字取自原文导言区，保证记号与原文完全一致。
%=====================================================================

%---------------------------------------------------------------------
% 一、页面与版式（A4，便于打印与屏幕阅读）
%---------------------------------------------------------------------
\usepackage[a4paper,left=2.6cm,right=2.6cm,top=2.6cm,bottom=2.6cm]{geometry}
\linespread{1.12}
\setlength{\emergencystretch}{3em}
\tolerance=2000
\setlength{\parindent}{2em}
\setlength{\parskip}{0.2ex}

%---------------------------------------------------------------------
% 二、中文支持（必须使用 XeLaTeX 内核）
%---------------------------------------------------------------------
\usepackage[UTF8,fontset=fandol]{ctex}
% 章标题格式：第 N 章  <标题>（居中）；小节标题保持 article 默认样式
\usepackage{titlesec}
\titleformat{\section}{\Large\bfseries\centering}{第\,\thesection\,章}{1em}{}
\titlespacing*{\section}{0pt}{3.2ex plus 1ex minus .2ex}{1.8ex plus .2ex}
\titleformat{\subsection}{\large\bfseries}{\thesubsection}{0.8em}{}
\titlespacing*{\subsection}{0pt}{2.4ex plus .6ex minus .2ex}{1ex plus .2ex}

%---------------------------------------------------------------------
% 三、数学宏包
%---------------------------------------------------------------------
\usepackage{amsmath}
\usepackage{amsthm}
\usepackage{amssymb}
\usepackage{amsfonts}
\usepackage{latexsym}
\usepackage{mathtools}

%---------------------------------------------------------------------
% 四、图形、颜色、表格、链接
%---------------------------------------------------------------------
\usepackage{graphicx}
\usepackage{xcolor}
\usepackage{booktabs}
\usepackage{array}
\usepackage{enumitem}
\setlist{nosep,leftmargin=*}

%---------------------------------------------------------------------
% 五、定理类环境（名称汉化，编号方式与原文完全一致：共用定理计数器、按章编号）
%---------------------------------------------------------------------
\theoremstyle{plain}
\newtheorem{theorem}{定理}[section]
\newtheorem{lemma}[theorem]{引理}
\newtheorem{exercise}[theorem]{习题}
\newtheorem{proposition}[theorem]{命题}
\newtheorem{corollary}[theorem]{推论}

\theoremstyle{definition}
\newtheorem{definition}[theorem]{定义}
\newtheorem{example}[theorem]{例}
\newtheorem{conjecture}[theorem]{猜想}

\theoremstyle{remark}
\newtheorem{remark}[theorem]{注}
\newtheorem{note}[theorem]{说明}

% 证明环境
\renewcommand{\proofname}{证明}

% 公式按章编号：(1.1), (1.2), ...
\numberwithin{equation}{section}

%---------------------------------------------------------------------
% 六、原文自带的全部记号缩写（逐字照搬，切勿改动）
%---------------------------------------------------------------------
% 本文件逐字取自原文导言区（qftlect.tex 第 23--133 行）中的宏定义，
% 保证所有记号、花体字、缩写与原文完全一致。

%%%%%%%%%%%%%%% symbols shorthand %%%%%%%%%%%  
%\renewcommand{\text}[1]{\mbox{{\rm #1}}}
\newcommand{\Tr}{\text{Tr}}%{\text{Tr}\,} 
\newcommand{\id}{\text{id}} 
\newcommand{\Image}{\text{Im}} 
\newcommand{\Ker}{\text{Ker\,}} 
\newcommand{\rank}{\text{rank\,}} 
\newcommand{\Index}{\text{Index}\,} 
\newcommand{\Map}{\text{Map}}
\newcommand{\End}{\text{End}} 
\newcommand{\PGL}{\text{PGL}}
\newcommand{\Hom}{\text{Hom}} 
\newcommand{\Ad}{\text{Ad}}
\newcommand{\Aut}{\text{Aut}}
\newcommand{\Sh}{\text{Sh}}
\newcommand{\Ind}{\text{Ind}}
\newcommand{\Rep}{\text{Rep}}
\newcommand{\supp}{\text{supp}\,} 
\renewcommand{\span}{\text{span}} 
\newcommand{\depth}{\text{depth}} 
\newcommand{\St}{\text{St}}
\newcommand{\diag}{\text{diag}}
\newcommand{\comp}{\text{comp}}
\newcommand{\op}{\text{op}}
\newcommand{\tr}{\text{tr}\,} 

   
\newcommand{\eps}{\varepsilon}
\newcommand{\mR}{{\mathcal R}} 
\newcommand{\bmR}{{\bar \mR}}
\newcommand{\bR}{{\bar R}}   
\newcommand{\bV}{{\bar V}}
\newcommand{\bF}{{\bar F}}
\newcommand{\I}{\mathcal{I}}
\renewcommand{\O}{\mathcal{O}}

\newcommand{\G}{\mathcal{G}}
\newcommand{\E}{\mathcal{E}}

\newcommand{\g}{\mathfrak{g}}
\newcommand{\h}{\mathfrak{h}}
\newcommand{\n}{\mathfrak{n}}
\renewcommand{\b}{\mathfrak{b}}
\newcommand{\wg}{\widehat{\mathfrak{g}}}
\newcommand{\wh}{\widehat{\mathfrak{h}}}
\newcommand{\wn}{\widehat{\mathfrak{n}}}
\newcommand{\wb}{\widehat{\mathfrak{b}}}
\newcommand{\ov}{\overline}

\newcommand{\actl}{\rightharpoonup}
\newcommand{\actr}{\leftharpoonup} 
\newcommand{\lact}{\triangleright} 
\newcommand{\ract}{\triangleleft} 
\newcommand{\la}{\langle\,} 
\newcommand{\ra}{\,\rangle} 
\newcommand{\flip}{\varsigma} 
\newcommand{\nv}{\vec{n}} 
\renewcommand{\star}{{\dag}} 
%\newcommand{\rtimes}{{>\!\!\!\triangleleft}} 
%\newcommand{\ltimes}{{\triangleright\!\!\!<}} 
\newcommand{\smrtimes}{>\!\!\triangleleft} 
\newcommand{\equivalent}{\Leftrightarrow} 
\newcommand{\half}{1/2}%%%%%{\frac{1}{2}}   
\newcommand{\own}{\ni}  
\newcommand{\1}{_{(1)}} 
\newcommand{\2}{_{(2)}} 
\newcommand{\3}{_{(3)}} 
\newcommand{\4}{_{(4)}}
\newcommand{\5}{_{(5)}} 
\newcommand{\6}{_{(6)}}
\renewcommand{\I}{^{(1)}} 
\newcommand{\II}{^{(2)}}   
\newcommand{\tDelta}{\Delta^{op}} 
\newcommand{\tS}{\widetilde{S}} 
\newcommand{\teps}{\widetilde\eps}  
\newcommand{\tI}{_{(\widetilde{1})}} 
\newcommand{\tII}{_{(\widetilde{2})}} 
\newcommand{\tIII}{_{(\widetilde{3})}}
\newcommand{\bTheta}{{\bar \Theta}}
\newcommand{\bs}{\backslash} 


\newcommand{\ot}{\otimes}
\newcommand{\qexp}{{\text{qexp}}}
\newcommand{\ben}{\begin{enumerate}}
\newcommand{\een}{\end{enumerate}}
\newcommand{\ad}{{\text{Ad}}}
\newcommand{\Rad}{{\text{Rad}}}
\newcommand{\Vect}{{\text{Vect}}}
\newcommand{\Lie}{{\text{Lie}}}
\newcommand{\Cl}{{\text{Cl}}}
\newcommand{\mC}{{\mathcal C}}
\newcommand{\trace}{{\rm tr}}
\newcommand{\CC}{{\mathbb{C}}}
\newcommand{\mM}{{\mathcal M}}

\hyphenation{se-mi-simple co-se-mi-simple}

%%%%%%% Bonnie added %%%%%%
\renewcommand{\Bbb}{\mathbb}
\renewcommand{\frak}{\mathfrak}
\renewcommand{\bold}{\mathbf}

%%%\newcommand\topfraction{.7}
\renewcommand\topfraction{.6}
%%%\newcommand\bottomfraction{.3}
\renewcommand\bottomfraction{.6}
\numberwithin{equation}{section}
\pagestyle{plain}

\makeindex


%---------------------------------------------------------------------
% 七、图表位置标注：在中文图题后自动注明原文页码与图号
%     用法：\caption{中文图题\figorig{原文页码}{原文图号}}
%---------------------------------------------------------------------
\newcommand{\figorig}[2]{\\[1pt]{\footnotesize\color{gray}（原文：第 #1 页，Figure #2）}}

%---------------------------------------------------------------------
% 八、每章“本节重点”环境（译者为方便阅读所加，非原文内容）
%---------------------------------------------------------------------
\newenvironment{keytakeaways}
  {\par\medskip\noindent
   \begin{minipage}{0.96\textwidth}
   \color{blue!45!black}\rule{\linewidth}{0.8pt}\par\smallskip
   \noindent\textbf{本章要点（Key Takeaways）}\par\smallskip\small}
  {\par\smallskip\noindent\color{blue!45!black}\rule{\linewidth}{0.8pt}\end{minipage}\medskip}

%---------------------------------------------------------------------
% 九、审校批注宏（正文一律用 % [Tutor Note: ...] 注释，此宏仅为备用）
%---------------------------------------------------------------------
\newcommand{\TutorNote}[1]{}

%---------------------------------------------------------------------
% 十、超链接（必须最后加载）
%---------------------------------------------------------------------
\usepackage[unicode,colorlinks=true,
            linkcolor=blue!55!black,
            citecolor=green!45!black,
            urlcolor=blue!60!black,
            pdftitle={量子场论的数学思想与概念},
            pdfauthor={Pavel Etingof（中译者）},
            bookmarksnumbered=true]{hyperref}

 读入。
%        宏定义 macros.tex 逐字取自原文导言区，保证记号与原文完全一致。
%=====================================================================

%---------------------------------------------------------------------
% 一、页面与版式（A4，便于打印与屏幕阅读）
%---------------------------------------------------------------------
\usepackage[a4paper,left=2.6cm,right=2.6cm,top=2.6cm,bottom=2.6cm]{geometry}
\linespread{1.12}
\setlength{\emergencystretch}{3em}
\tolerance=2000
\setlength{\parindent}{2em}
\setlength{\parskip}{0.2ex}

%---------------------------------------------------------------------
% 二、中文支持（必须使用 XeLaTeX 内核）
%---------------------------------------------------------------------
\usepackage[UTF8,fontset=fandol]{ctex}
% 章标题格式：第 N 章  <标题>（居中）；小节标题保持 article 默认样式
\usepackage{titlesec}
\titleformat{\section}{\Large\bfseries\centering}{第\,\thesection\,章}{1em}{}
\titlespacing*{\section}{0pt}{3.2ex plus 1ex minus .2ex}{1.8ex plus .2ex}
\titleformat{\subsection}{\large\bfseries}{\thesubsection}{0.8em}{}
\titlespacing*{\subsection}{0pt}{2.4ex plus .6ex minus .2ex}{1ex plus .2ex}

%---------------------------------------------------------------------
% 三、数学宏包
%---------------------------------------------------------------------
\usepackage{amsmath}
\usepackage{amsthm}
\usepackage{amssymb}
\usepackage{amsfonts}
\usepackage{latexsym}
\usepackage{mathtools}

%---------------------------------------------------------------------
% 四、图形、颜色、表格、链接
%---------------------------------------------------------------------
\usepackage{graphicx}
\usepackage{xcolor}
\usepackage{booktabs}
\usepackage{array}
\usepackage{enumitem}
\setlist{nosep,leftmargin=*}

%---------------------------------------------------------------------
% 五、定理类环境（名称汉化，编号方式与原文完全一致：共用定理计数器、按章编号）
%---------------------------------------------------------------------
\theoremstyle{plain}
\newtheorem{theorem}{定理}[section]
\newtheorem{lemma}[theorem]{引理}
\newtheorem{exercise}[theorem]{习题}
\newtheorem{proposition}[theorem]{命题}
\newtheorem{corollary}[theorem]{推论}

\theoremstyle{definition}
\newtheorem{definition}[theorem]{定义}
\newtheorem{example}[theorem]{例}
\newtheorem{conjecture}[theorem]{猜想}

\theoremstyle{remark}
\newtheorem{remark}[theorem]{注}
\newtheorem{note}[theorem]{说明}

% 证明环境
\renewcommand{\proofname}{证明}

% 公式按章编号：(1.1), (1.2), ...
\numberwithin{equation}{section}

%---------------------------------------------------------------------
% 六、原文自带的全部记号缩写（逐字照搬，切勿改动）
%---------------------------------------------------------------------
% 本文件逐字取自原文导言区（qftlect.tex 第 23--133 行）中的宏定义，
% 保证所有记号、花体字、缩写与原文完全一致。

%%%%%%%%%%%%%%% symbols shorthand %%%%%%%%%%%  
%\renewcommand{\text}[1]{\mbox{{\rm #1}}}
\newcommand{\Tr}{\text{Tr}}%{\text{Tr}\,} 
\newcommand{\id}{\text{id}} 
\newcommand{\Image}{\text{Im}} 
\newcommand{\Ker}{\text{Ker\,}} 
\newcommand{\rank}{\text{rank\,}} 
\newcommand{\Index}{\text{Index}\,} 
\newcommand{\Map}{\text{Map}}
\newcommand{\End}{\text{End}} 
\newcommand{\PGL}{\text{PGL}}
\newcommand{\Hom}{\text{Hom}} 
\newcommand{\Ad}{\text{Ad}}
\newcommand{\Aut}{\text{Aut}}
\newcommand{\Sh}{\text{Sh}}
\newcommand{\Ind}{\text{Ind}}
\newcommand{\Rep}{\text{Rep}}
\newcommand{\supp}{\text{supp}\,} 
\renewcommand{\span}{\text{span}} 
\newcommand{\depth}{\text{depth}} 
\newcommand{\St}{\text{St}}
\newcommand{\diag}{\text{diag}}
\newcommand{\comp}{\text{comp}}
\newcommand{\op}{\text{op}}
\newcommand{\tr}{\text{tr}\,} 

   
\newcommand{\eps}{\varepsilon}
\newcommand{\mR}{{\mathcal R}} 
\newcommand{\bmR}{{\bar \mR}}
\newcommand{\bR}{{\bar R}}   
\newcommand{\bV}{{\bar V}}
\newcommand{\bF}{{\bar F}}
\newcommand{\I}{\mathcal{I}}
\renewcommand{\O}{\mathcal{O}}

\newcommand{\G}{\mathcal{G}}
\newcommand{\E}{\mathcal{E}}

\newcommand{\g}{\mathfrak{g}}
\newcommand{\h}{\mathfrak{h}}
\newcommand{\n}{\mathfrak{n}}
\renewcommand{\b}{\mathfrak{b}}
\newcommand{\wg}{\widehat{\mathfrak{g}}}
\newcommand{\wh}{\widehat{\mathfrak{h}}}
\newcommand{\wn}{\widehat{\mathfrak{n}}}
\newcommand{\wb}{\widehat{\mathfrak{b}}}
\newcommand{\ov}{\overline}

\newcommand{\actl}{\rightharpoonup}
\newcommand{\actr}{\leftharpoonup} 
\newcommand{\lact}{\triangleright} 
\newcommand{\ract}{\triangleleft} 
\newcommand{\la}{\langle\,} 
\newcommand{\ra}{\,\rangle} 
\newcommand{\flip}{\varsigma} 
\newcommand{\nv}{\vec{n}} 
\renewcommand{\star}{{\dag}} 
%\newcommand{\rtimes}{{>\!\!\!\triangleleft}} 
%\newcommand{\ltimes}{{\triangleright\!\!\!<}} 
\newcommand{\smrtimes}{>\!\!\triangleleft} 
\newcommand{\equivalent}{\Leftrightarrow} 
\newcommand{\half}{1/2}%%%%%{\frac{1}{2}}   
\newcommand{\own}{\ni}  
\newcommand{\1}{_{(1)}} 
\newcommand{\2}{_{(2)}} 
\newcommand{\3}{_{(3)}} 
\newcommand{\4}{_{(4)}}
\newcommand{\5}{_{(5)}} 
\newcommand{\6}{_{(6)}}
\renewcommand{\I}{^{(1)}} 
\newcommand{\II}{^{(2)}}   
\newcommand{\tDelta}{\Delta^{op}} 
\newcommand{\tS}{\widetilde{S}} 
\newcommand{\teps}{\widetilde\eps}  
\newcommand{\tI}{_{(\widetilde{1})}} 
\newcommand{\tII}{_{(\widetilde{2})}} 
\newcommand{\tIII}{_{(\widetilde{3})}}
\newcommand{\bTheta}{{\bar \Theta}}
\newcommand{\bs}{\backslash} 


\newcommand{\ot}{\otimes}
\newcommand{\qexp}{{\text{qexp}}}
\newcommand{\ben}{\begin{enumerate}}
\newcommand{\een}{\end{enumerate}}
\newcommand{\ad}{{\text{Ad}}}
\newcommand{\Rad}{{\text{Rad}}}
\newcommand{\Vect}{{\text{Vect}}}
\newcommand{\Lie}{{\text{Lie}}}
\newcommand{\Cl}{{\text{Cl}}}
\newcommand{\mC}{{\mathcal C}}
\newcommand{\trace}{{\rm tr}}
\newcommand{\CC}{{\mathbb{C}}}
\newcommand{\mM}{{\mathcal M}}

\hyphenation{se-mi-simple co-se-mi-simple}

%%%%%%% Bonnie added %%%%%%
\renewcommand{\Bbb}{\mathbb}
\renewcommand{\frak}{\mathfrak}
\renewcommand{\bold}{\mathbf}

%%%\newcommand\topfraction{.7}
\renewcommand\topfraction{.6}
%%%\newcommand\bottomfraction{.3}
\renewcommand\bottomfraction{.6}
\numberwithin{equation}{section}
\pagestyle{plain}

\makeindex


%---------------------------------------------------------------------
% 七、图表位置标注：在中文图题后自动注明原文页码与图号
%     用法：\caption{中文图题\figorig{原文页码}{原文图号}}
%---------------------------------------------------------------------
\newcommand{\figorig}[2]{\\[1pt]{\footnotesize\color{gray}（原文：第 #1 页，Figure #2）}}

%---------------------------------------------------------------------
% 八、每章“本节重点”环境（译者为方便阅读所加，非原文内容）
%---------------------------------------------------------------------
\newenvironment{keytakeaways}
  {\par\medskip\noindent
   \begin{minipage}{0.96\textwidth}
   \color{blue!45!black}\rule{\linewidth}{0.8pt}\par\smallskip
   \noindent\textbf{本章要点（Key Takeaways）}\par\smallskip\small}
  {\par\smallskip\noindent\color{blue!45!black}\rule{\linewidth}{0.8pt}\end{minipage}\medskip}

%---------------------------------------------------------------------
% 九、审校批注宏（正文一律用 % [Tutor Note: ...] 注释，此宏仅为备用）
%---------------------------------------------------------------------
\newcommand{\TutorNote}[1]{}

%---------------------------------------------------------------------
% 十、超链接（必须最后加载）
%---------------------------------------------------------------------
\usepackage[unicode,colorlinks=true,
            linkcolor=blue!55!black,
            citecolor=green!45!black,
            urlcolor=blue!60!black,
            pdftitle={量子场论的数学思想与概念},
            pdfauthor={Pavel Etingof（中译者）},
            bookmarksnumbered=true]{hyperref}

